\chapter{Trabalhos Relacionados}
\label{cap:trabalhos-relacionados}

\citeonline{Lewis2020} combinaram um recuperador denso com um gerador
\textit{sequence-to-sequence} e mostraram que condicionar a geração em passagens
recuperadas melhora tarefas que dependem de conhecimento factual. O levantamento de
\citeonline{Gao2023} organiza as variações posteriores em RAG ingênuo, avançado e
modular, e identifica a divisão dos documentos em trechos e a escolha do modelo de
\textit{embedding} como pontos que afetam diretamente a qualidade da recuperação. Este
trabalho fica na categoria mais simples, recuperar e depois gerar, e concentra o
esforço na divisão em trechos, que segue a estrutura das normas.

Para textos em português, a escolha do \textit{embedding} pesa. O BGE-M3
\cite{Chen2024} foi treinado para mais de cem línguas e aceita entradas longas, o que o
torna adequado a artigos extensos de regimentos e resoluções.

A avaliação de sistemas RAG costuma separar a recuperação da geração. O RAGAS
\cite{Es2023} propõe métricas sem referência humana, como fidelidade da resposta ao
contexto recuperado, relevância da resposta e precisão do contexto, calculadas com
auxílio de um LLM. \citeonline{Zheng2023} estudaram o uso de LLMs como juízes e
encontraram boa concordância com avaliadores humanos, mas também vieses de posição, de
verbosidade e de autopreferência. Esses vieses motivam, neste trabalho, o uso de um
juiz de família diferente da dos modelos avaliados e a comparação explícita entre as
notas do juiz e as notas humanas.

Do lado da segurança, \citeonline{Greshake2023} mostraram que aplicações integradas a
LLMs podem ser comprometidas por injeção indireta de \textit{prompt}: instruções
escondidas em conteúdo recuperado, e não digitadas pelo usuário, passam a controlar o
modelo. Em um sistema RAG qualquer documento indexado é uma superfície de ataque, e
esse cenário é testado aqui com um documento envenenado. A alucinação, discutida em
profundidade por \citeonline{Ji2023}, é medida pela taxa de respostas sem fonte e pela
nota de fidelidade.

O ajuste fino eficiente com LoRA \cite{Hu2021} reduz o custo de adaptar um modelo,
mas incorpora o conhecimento aos pesos. Para normas que mudam a cada resolução, isso
exigiria retreinar o modelo a cada alteração e não resolveria a necessidade de citar
a fonte, razão pela qual a alternativa foi descartada (Seção~\ref{sec:estrategia}).
